\section{Introduction}\label{sec:intro}

\paragraph{The question.} Can a learner without continuous feature parameters do at least as well as gradient-descent-trained networks, with the same training data and no more FLOPs? This paper reports the latest learner of a study that has asked this question in ten parts, and what that learner does.

\paragraph{The study so far.} Each part answered a question posed by the user who prompted the work. From Part 6 on, each part froze its learners before running them on the held-out data that it reports as its main evidence.
\begin{itemize}
\item \emph{Parts 1--5} \cite{gdcircuit2026} trained rich-regime networks on Boolean targets and described what gradient descent does there. It grows features greedily. A feature's cost is set by how hard its correlation with the current residual is to detect against the sampling noise of the competing directions, and features on inputs already in use are cheaper. The price is paid in samples. For a $k$-bit parity on $d$ inputs, the networks' sample size at 10\% test error, $\nstar$, grows like $d^{1.0}$, $d^{1.9}$ and $d^{2.4}$ for $k=2,3,4$. That is close to the leap and correlational-SQ rate $d^{k-1}$ \cite{abbe2022msp,abbe2023leap} and far from the $\approx k\log d$ of an Occam learner.
\item \emph{Part 6} built DFG (discrete feature growth) from that picture \cite{discretelearners2026}. DFG enumerates Fourier characters and admits each by a union-bound test against the residual. On a blind 7-target suite it needed less data than the networks on every target (median factor 3.1), with at most 0.26 of one network run's FLOPs per fit.
\item \emph{Part 7} ran three adaptations of DFG on standard tasks. They were about level with a plain MLP on tabular data, behind a regularized one, and clearly behind on images at large $n$.
\item \emph{Part 8} stacked DFG with LightGBM, ExtraTrees, nearest-neighbour votes and logistic regression (v1). v1 needed clearly less data than the best network variant on 9 of 10 fresh toy targets. Its mean error was lower than the hindsight-best MLP's on 8 of 10 fresh tabular datasets (sign test $p=0.11$). On images it was behind a regularized MLP ensemble at $n\ge30000$.
\item \emph{Part 9} added sum blocks to DFG and permutation-invariant image members to the stack (v2). On a 12-target suite written blind to the sum-feature design, v2 needed clearly less data than the best pre-registered network variant on 10 targets and more on none. On two fresh QuickDraw sets it beat a strong MLP ensemble at $n\ge10000$. But a CPU-matched ensemble of wider MLPs and an MLP given v2's own pooled features tied it at $n=60000$, and a small CNN that sees the pixel layout was clearly better at $n=10000$.
\end{itemize}

\paragraph{Part 10, and what this paper adds.} The user asked for ``a couple more hours'' on a better discrete learner, then for interpretability of the best learner so far, and then for a paper about it that explains what kind of learner it is, how well it does against the baselines, and what the interpretability shows. Part 10 pre-registered the whole sequence (\cref{sec:discipline}). Three independent builders worked on development data only. A rule fixed in advance chose two of their components, a new sum-feature search for DFG and a second image booster on ``bar'' features. The assembled learner v3 was frozen before any of its held-out results were read, and a second rule fixed in advance made v3, not v2, the subject of this paper. The paper contributes four things:
\begin{enumerate}
\item an explanation of what v3 is and what in it is discrete (\cref{sec:learner});
\item a held-out test of v3 on a fresh toy suite and two fresh image sets, against the pre-registered references (\cref{sec:res-dlheld7,sec:res-img-heldout});
\item a comparison of v3 with every reference on every toy suite and real dataset of the study, with the status of each set for v3 stated (\cref{sec:results-toy,sec:results-img});
\item a descriptive analysis of what v3's fitted members contain, checked by independent verifiers (\cref{sec:interp}).
\end{enumerate}

\paragraph{What was found.}
\begin{itemize}
\item \textbf{Held-out toy suite} \status{Numerical} (\cref{sec:res-dlheld7}). The suite is \code{dlheld7}, 12 targets drawn from the blind family of Part 9. v3's mean $\log_2\nstar$ ratio against the per-target best of three pre-registered network variants is $-1.50$ over the 11 targets both learn, about $2.8\times$ less data. v3 is clearly better on 10 targets and clearly worse on none. One target is a tie (a halfspace), and one is learned by no learner (a residue class mod 3 of a 14-input count). On five of the ten wins, default gradient boosting or logistic regression already beats the networks, and there v3 leads them by small factors only. The Part 10 toy component changes nothing measurable on this suite: against v2 the mean is $-0.004$.
\item \textbf{Every toy suite} \status{Numerical} (\cref{sec:res-toy-all}). On the 8 suites with v3 and network results so far, v3's suite means against the per-suite best network range from $-1.50$ to $-2.44$, and a network is clearly better on one target only (h6\_00, a residue class mod 3 that only the most regularized network learns, at $\nstar=12304$). Against the one network variant run on every suite, the margin is smallest on the two draws of the blind family. These suites are development data for v3, except \code{dlheld7}.
\item \textbf{Held-out image sets} \status{Numerical} (\cref{sec:res-img-heldout}). On \code{qd_c} and \code{qd_d}, v3's test error at $n=60000$ is 12.76\% and 12.64\%. It is clearly better than the three-network MLP ensemble \code{mlp_strong} at all six cells (errors 17--23\% lower), than gradient boosting and v1 at all six, and than \code{mlp_pool}, an MLP given v2's pooled features, at both cells where that ran. It is clearly better than the 512-wide ensemble at five cells and ties it at the sixth. It ties the CPU-matched 1024-wide ensemble at both of its cells (8.6\% and 9.6\% lower error). It ties v2 at all six cells, with lower error at each (2.2--9.6\%).
\item \textbf{Tabular data} \status{Numerical} (\cref{sec:res-tabular}). On 24 of the 26 datasets v3 runs v1's code, unchanged since Part 8. On the 10 datasets fixed in Part 8 its mean $\log_2$ error ratio against the hindsight-best MLP is $-0.168$, lower on 8 of 10 datasets ($p=0.11$).
\item \textbf{What DFG learns} \status{Numerical} (\cref{sec:interp-toy}). On 49 of 54 sum-structured targets, DFG's model comes to hold the target's exact signed blocks, at a median $n$ of 256. On the 25 targets that are functions of one signed sum, DFG's test would admit the true block at $n=128$ every time. The extra data needed to find it grows as the target's Fourier weight at degrees 1--2 shrinks (Spearman $\rho=-0.86$). The three residue classes mod 3 have no weight there, and the search never proposes their block, although at $n=16384$ it would pass the test at 134--155 times the threshold.
\item \textbf{What the image members encode} \status{Numerical} (\cref{sec:interp-img}). The pooling sets are local neighbourhoods of the true pixel grid on QuickDraw and MNIST (98.3--99.9\% of the other members of a 4-pixel set lie within one step on QuickDraw), and mirror-image pairs at some pixels of Fashion-MNIST. The bar booster that v3 adds has the lowest out-of-fold error of all members in every analysed fit and takes over the old booster's stacking weight.
\item \textbf{What kind of learner it is} \status{Heuristic} (\cref{sec:discussion}). v3 is three learners under one stacker, chosen by input type. On Boolean targets it is enumerate-and-test, at about an Occam sample cost wherever its search can propose the right feature. On images it is a fixed hierarchy of data-derived subset sums read out by gradient-boosted trees. On other tabular data it is a classical stack.
\end{itemize}

\paragraph{Scope.}
\begin{itemize}
\item \emph{What is held out.} For v3 as a whole only three sets are held out: \code{dlheld7}, \code{qd_c} and \code{qd_d} (\cref{tab:status}). They were fixed in a commit before any Part 10 design work, and no result of a Part 10 learner on them was read before v3 was frozen. The earlier held-out suites and image sets were development data in Part 10 or shaped earlier versions, so they measure v3 but do not test it. On tabular data the cleaner claim is about v1's code path, which v3 runs there unchanged.
\item \emph{The references.} They are concrete networks with standard recipes. On the toy suites they are $\mu$P two-hidden-layer MLPs of width 128 trained by Adam, in up to five loss and weight-decay settings. On images they are two-hidden-layer MLPs of width 256--1024, up to ensembles of three with weight decay and dropout. There is no augmentation and no architecture search. The per-target and per-cell minima \code{nets_best} and \code{mlp_best} are hindsight choices that favour the references. A small CNN on the unpermuted pixels is reported as context, not as a reference.
\item \emph{What ``discrete'' means.} No member fits a continuous feature direction by gradient descent. Features are enumerated, built from data statistics such as correlation ranks, and admitted by tests or by greedy split search. The continuous fits are linear readouts, tree leaf values and stacking weights. Gradient boosting is functional gradient descent over trees, and two components are borderline (\cref{sec:whatdiscrete}). The fair wording is ``compared with gradient-descent-trained networks'', not ``compared with gradient descent''.
\item \emph{Power.} The held-out sets are 12 toy targets from one family and two image sets from one source, with one seed at $n=60000$.
\end{itemize}
Claims are labelled \status{Numerical} (measured here, with the source file named) or \status{Heuristic} (a reading or a back-of-envelope argument, not tested here). There are no theorems in this paper.

\paragraph{Organisation.} \Cref{sec:learner} explains the learner and \cref{sec:protocol} the tasks, references, metrics and the pre-registration. \Cref{sec:results-toy} gives the toy and tabular results and \cref{sec:results-img} the image results. \Cref{sec:interp} reports the interpretability analyses, and \cref{sec:discussion} discusses what the results mean and where they stop. \Cref{app:provenance} records how Part 10 was run.
